\documentclass[11pt,a4paper]{article}
\usepackage[margin=23mm]{geometry}
\usepackage{fontspec}\setmainfont{Latin Modern Roman}
\usepackage{amsmath,amssymb,amsthm,xurl,hyperref}
\setlength{\parindent}{0pt}\setlength{\parskip}{6pt}
\setlength{\emergencystretch}{3em}
\newcommand{\code}[1]{\nolinkurl{#1}}
\newcommand{\none}{\mathtt{none}}
\newcommand{\E}{\mathbb E}
\newtheorem{theorem}{Twierdzenie}
\title{S05 / MARGINAL\_006\\\large Ważony marginal, capy i nowy certyfikat pełnej odpowiedzi}
\author{Notatka robocza dla Niirmata}\date{10 października 2026}
\begin{document}\maketitle
\textbf{WARUNKOWY / DEVELOPMENT / textual.}
Nowa konstrukcja otrzymuje własny dowód; nie dziedziczy certyfikatu h=0.
Wynik to warunkowy interfejs certyfikatu z małym $e$, dokładne skończone
kontrole prawa oraz przeszkoda kosztowa dla konkretnej propozycji.
Nie skonstruowano efektywnego świadka Q-JOINT-INSTANCE dla wszystkich $h$.

\section{Właściwy marginal i dane do jego realizacji}
Zachowujemy $D,w,G,Z_1,Z_{h,c}$ z NORMALIZER-005, $q=18433$, $\sigma=768$,
768 bloków w połówce oraz 16 prób uczciwego testu normy.
Dla ustalonych $h,c$ piszmy $F(y)=Z_1(c-hy)$ i $U(y)=w(y)/G$.
Właściwy marginal drugiej połówki to dokładnie
\[
 V(y)=\frac{w(y)F(y)}{Z_{h,c}}=\frac{U(y)F(y)}{\E_U F}.
\]
Normalizator globalny znika z implementacji odrzucania, ale pozostaje
w prawdopodobieństwie jego akceptacji.
NORMALIZER-005 daje $F_-(y)\le F(y)\le\kappa F_-(y)$,
$\kappa=(1+2^{-366})^{768}<1+2^{-356}$, dla każdego $y$.
Ta kontrola względna jest dokładnie kontrolą wag potrzebną marginalowi.
Numeryczne wejście pochodzi ze skorygowanego certyfikatu 005:
\code{NORMALIZER_005_ERRATUM_001.tex} opisuje wykryty błąd
\code{QQ(MPFR)} i ponowny rachunek z \code{exact_rational()}.

Wybierzmy stałą $M\ge\max_yF(y)$, rozkład propozycji $\nu$ i monetę
akceptacji $a(y)$. Wymagamy jawnych certyfikatów
\[
 \ell_U U\le\nu\le u_U U,\qquad
 \ell_F F/M\le a\le u_F F/M,\qquad \ell_U,\ell_F>0.
\]
W szczególności wszystkie istotne nośniki pozostają obecne.
Po akceptacji pobieramy $z_1$ z przybliżonego prawa kategorii w reszcie
$c-hy$, z dominacją $K_1$ względem dokładnego Gaussa warunkowego.
\textbf{Propozycja 004 nie spełnia tych założeń na mocy swojego certyfikatu:}
ucina współrzędne do 15360 i ma jednostronny certyfikat z fallbackiem.
Tutaj potrzebny jest osobny certyfikat propozycji na pełnym boxie albo
odrębne oszacowanie jej znormalizowanej masy po ważeniu.

\subsection*{Względna dokładność skończonej monety}
Jeżeli wykonalny dolny przedział $L(y)$ spełnia
$F(y)/\kappa\le L(y)\le F(y)$, to
$a(y)=2^{-b}\lfloor2^b L(y)/M\rfloor$ spełnia
\[
 (1-\eta)F(y)/(\kappa M)\le a(y)\le F(y)/M
\]
pod warunkiem $2^{-b}\le\eta\min_y L(y)/M$.
Nie wystarcza mały błąd \emph{bezwzględny} monety, gdy $F(y)/M$ jest małe.
Symboliczne końce z 005 zawierają czynnik wykładniczy; realizacja $L(y)$
przez skończone operacje i rozstrzyganie progów wymaga własnego rachunku
zaokrągleń. Nie utożsamiamy zapisu $F_-$ z już wykonanym generatorem monety.

\section{Certyfikat po akceptacji, wyprowadzony od początku}
Prawdopodobieństwo akceptacji jednej propozycji jest
$\alpha=\sum_y\nu(y)a(y)>0$, a marginal po akceptacji to
$\widehat V(y)=\nu(y)a(y)/\alpha$.
Ponieważ $\alpha\ge\ell_U\ell_F\E_U F/M$,
\[
 \widehat V(y)\le\frac{u_Uu_F}{\ell_U\ell_F}V(y).
\]
Po kategorii $z_1$ jedna zaakceptowana para ma prawo $Q_A$ spełniające
\[
 Q_A\le dP_A,\qquad d=\frac{u_Uu_F}{\ell_U\ell_F}K_1,
\]
gdzie $P_A$ to uczciwe prawo całego włókna przed normą.
Po 16 niezależnych parach i tym samym kanale normy, cap16 oraz emisji
otrzymujemy pełne prawo $Q$ z
\[
 Q\le CP,\quad C=d^{16},\quad Q\ll P,\quad
 \sum Q^2/P\le C,\quad\chi^2(Q\Vert P)\le C-1.
\]
Ostatnia granica wynika z $Q^2/P\le CQ$ i $\sum Q=1$.
Obejmuje $\none$; nie jest to argument TV zamieniony na moment.

\section{Dokładne prawo nowego capu}
\textbf{Definicja konstrukcji:} przygotuj wszystkie 16 par. Dla każdej
wolno wykonać najwyżej $N$ niezależnych propozycji marginalu.
Jeżeli jakikolwiek wewnętrzny cap się wyczerpie, wydaj $\none$.
W przeciwnym przypadku zastosuj do 16 par uczciwy test normy, wybór
pierwszej akceptacji i modelową emisję. Maksymalnie jest $16N$ propozycji.

Prawo zaakceptowanej pary nie zależy od numeru akceptacji:
masa akceptacji $y$ w kroku $k$ to $(1-\alpha)^{k-1}\nu(y)a(y)$.
Po zsumowaniu do $N$ i normalizacji otrzymujemy to samo $\widehat V$.
Z niezależności 16 bloków pracy wynika więc dokładnie
\[
 \delta=1-\bigl(1-(1-\alpha)^N\bigr)^{16},\qquad
 J=(1-\delta)Q+\delta\,\delta_{\none}.
\]
Nie zastępujemy tej konstrukcji leniwym zatrzymaniem po pierwszym sukcesie
normy: taki algorytm wymagałby nowego rozliczenia swojego prawa.
Koszt i gałąź $\none$ są częścią definicji nowego kandydata.

\begin{theorem}[Nowy certyfikat pełnego prawa z capami]
Załóżmy powyższe warunki i $p=P(\none)>0$. Wtedy $J\ll P$ oraz
\begin{align*}
 \sum J^2/P
 &=(1-\delta)^2\sum Q^2/P
   +2(1-\delta)\delta\,Q(\none)/p+\delta^2/p,\\
 \chi^2(J\Vert P)
 &\le\bigl((1-\delta)\sqrt{C-1}+\delta\sqrt{1/p-1}\bigr)^2\\
 &\le2(1-\delta)^2(C-1)+2\delta^2(1/p-1).
\end{align*}
Dla dokładnego rdzenia $Q=P$ mamy równość
$\sum J^2/P=1+\delta^2(1/p-1)$.
\end{theorem}
Dowód równości to rozwinięcie kwadratu, z osobnym składnikiem dla $\none$.
Dla nierówności piszemy w $L^2(P)$
$J/P-1=(1-\delta)(Q/P-1)+\delta(\mathbf1_{\none}/p-1)$,
stosujemy nierówność trójkąta, a następnie $(a+b)^2\le2a^2+2b^2$.
Nośnik zachowujemy, bo obie składowe mieszanki są absolutnie ciągłe względem $P$.
Jeżeli $p=0$ i $\delta>0$, dodanie tej porażki \emph{łamie} $J\ll P$.
Nie wolno więc traktować wewnętrznego abortu jako darmowej porażki.

Przy dolnych certyfikatach $\alpha\ge\alpha_0>0$, $p\ge\beta_0>0$,
\[
 \delta\le16e^{-N\alpha_0},\qquad
 \chi^2(J\Vert P)\le2(C-1)+512e^{-2N\alpha_0}/\beta_0.
\]
To warunek wyboru capu, a nie test jego poprawności przez kilka uruchomień.

\section{Warunkowy budżet z użytecznym e}
Weźmy $\rho=\eta=2^{-256}$, $\ell_U=1-\rho$, $u_U=1+\rho$,
$\ell_F=(1-\eta)/\kappa$, $u_F=1$ oraz
\[
 K_1=\bigl(1+64^2(2^{-372}+2^{-256})\bigr)^{768}.
\]
Dokładny rachunek $\mathbb Q$ daje $C-1<2^{-230}$.
Jeśli dodatkowo $\delta^2/\beta_0\le2^{-232}$, to
$\chi^2(J\Vert P)<2^{-228}$ dla każdego $h,c$ objętego tymi przesłankami.

Losując wyzwanie przez 1536 niezależnych 80-bitowych słów modulo $q$,
otrzymujemy dokładny moment wyzwania
\[
 M_H=\left(1+\frac{r(q-r)}{2^{160}}\right)^{1536},\qquad r=2^{80}\bmod q.
\]
Zsumowanie pełnego wspólnego prawa po $c$ daje
\[
 J_{\rm joint}\ll P_{\rm joint},\qquad
 \sum J_{\rm joint}^2/P_{\rm joint}
 \le M_H(1+2^{-228})<1+2^{-119}.
\]
Rachunek nie używa odziedziczonego wniosku $e<2^{-50}$ z h=0.
\textbf{Jest warunkowy:} propozycja, moneta, dodatnie masy i koszt są
nowymi obowiązkami; nie odhaczono Q-JOINT-INSTANCE.

\section{Dodatnia masa porażki istnieje, lecz prosty cap jest ogromny}
W skończonym modelu boxu dla każdego $h,c$ można wybrać $z_2$ mające
jeden blok $(H,H)$ i zera poza nim, $H=65535$. Jego norma już przekracza
$B=2093922385$. Dobieramy centrowane lifty $z_1$ do $c-hz_2$;
ich energia jest co najwyżej $E_1=768\cdot3\cdot9216^2$.
Ponieważ normalizator włókna jest mniejszy niż $2^{17\cdot3072}$,
a $e<4$, prawdopodobieństwo odrzucenia normy jednej próby jest większe niż
\[
 2^{-R_0},\qquad
 R_0=2\left\lceil\frac{E_1+3H^2}{2\sigma^2}\right\rceil+17\cdot3072.
\]
Stąd $p\ge\beta_0=2^{-16R_0}$. Nie zakładamy porażki emisji, aby to uzyskać.
To jawna, bardzo luźna dolna granica. Nie dowodzi typowej częstości porażek.

Dla $M=64^{768}$ mamy $F/M>2^{-336384}$, bo
$E_1/(2\sigma^2)=165888$ i znów $e<4$.
Przy zadanych względnych błędach daje to $\alpha_0=2^{-336385}$.
Wystarczy $N\alpha_0\ge120\log2+\tfrac12\log(1/\beta_0)$, aby
$\delta^2/\beta_0\le2^{-232}$.
Wykładnik wybranego dostatecznego $N$ zapisuje certyfikat; takiego capu
\emph{nie wykonywano}. Skończoność nie oznacza wykonalnego kosztu reduktora.
Dostateczny budżet monety $b=336641$ bitów obejmuje osobno podłogę dyadyczną,
zakładając dostęp do certyfikowanego $L$. Nie rozwiązuje generowania $L$.

\section{Przeszkoda kosztowa dla odrzucania całego wektora}
Ta przeszkoda dotyczy nawet \emph{idealnej} propozycji $U=w/G$ i dokładnej
monety $F/M$. Dla $h=1,c=0$ mamy $M\ge F(0)\ge1$ oraz
\[
 \alpha=Z_{1,0}/(GM)
 \le\left(Z_{\rm block}/G_\triangle\right)^{768}<2^{-767}.
\]
Górny certyfikat $Z_{\rm block}$ pochodzi z NORMALIZER-005.
$G_\triangle$ obliczono w balls, przez sumy theta, ogony geometryczne
i osobny koszt przejścia z nieskończonej kraty do boxu.
Wobec tego nawet po $N=2^{128}$ propozycjach
\[
 \Pr[\text{co najmniej jedna akceptacja}]\le N\alpha<2^{-639}.
\]
To dolna granica trudności \emph{tej strategii odrzucania całego wektora},
nie dowód niemożliwości innego marginalnego samplera.
Dla strukturalnego h=1 faktoryzacja blokowa omija ten konkretny iloczyn
akceptacji; dla ogólnego $h$ potrzebna jest adaptowana propozycja lub inna
kontrolowana konstrukcja. Nie zamieniamy małej akceptacji w twierdzenie
o dużym końcowym momencie po wyzwaniu i cap16.

\section{Wykonanie i następne obowiązki}
\code{MARGINAL_006.sage} sprawdza 102 dokładne prawa skończone
($q=3,5$, wszystkie $h,c$, capy 1,3,17), w tym emisję, normę, cap16,
pełną mieszankę abortu, równość momentu i dominację od założeń.
Osobna rekurencja po stanach 16 przygotowań, w której późny wewnętrzny
abort nadpisuje nawet już wybrany sukces, potwierdza prawo mieszanki
bez używania jej wzoru w tej rekurencji.
Osobna kontrola negatywna pokazuje utratę nośnika przy $P(\none)=0$.
Te kontrole są testami rachunku prawa; ogólny dowód jest tekstowy powyżej.
Rachunek FT1536 wykonano przy 512 i 768 bitach, z dokładnymi końcami przedziałów.

Otwarte: pełnoboxowa propozycja ze świeżym certyfikatem, skończona moneta
o względnej dokładności, użyteczny koszt marginalu dla wszystkich $h,c$,
kernelizacja h=0 i nowego lematu capów oraz źródłowe i kwantowe mosty.
Obie wcześniejsze przeszkody i nowy wynik kosztowy są zebrane w
\code{obstructions/}. Jest to materiał roboczy do ewentualnego Aneksu B;
paper, strona i statusy T12.1/B20 pozostają bez zmian.
\end{document}
